\subsection{The Fourth Sunday}\label{sec:fourth-sunday}

The last Sunday of Advent opens with Isaiah's cry for the Just One. Its Entrance
Antiphon is, word for word, the Introit the Roman Missal sang before the reform,
taken from Isaiah:

\begin{quote}
\latin{Roráte, cæli, désuper, et nubes pluant justum: aperiátur terra, et
gérminet Salvatórem.}\par
{\smallskip\noindent\small ---Introit of the Fourth Sunday of Advent\par}
\end{quote}

\begin{quote}
Drop down dew, ye heavens, from above, and let the clouds rain the just; let the
earth be opened and bud forth a Saviour.\par
{\smallskip\noindent\small ---The same, in English\par}
\end{quote}

\noindent The Gospel of this Sunday comes to the threshold of the birth. In Year
B it is the Annunciation. The first reading is the promise God
made to David through the prophet Nathan, when the king wished to build the Lord
a house and the Lord promised instead to build David's:

\begin{quote}
And when thy days shall be fulfilled, and thou shalt sleep with thy fathers, I
will raise up thy seed after thee, which shall proceed out of thy bowels, and I
will establish his kingdom. He shall build a house to my name, and I will
establish the throne of his kingdom for ever. \ldots{} And thy house shall be
faithful, and thy kingdom for ever before thy face, and thy throne shall be firm
for ever.\par
{\smallskip\noindent\small ---2 Samuel 7:12--13, 16 (2 Kings in the Douay)\par}
\end{quote}

\noindent St Augustine, in the \emph{City of God}, asked what house the son of
David would build:

\begin{quote}
For however long the interval of time might be before Jesus Christ came, beyond
doubt it was after the death of king David, to whom He was so promised, that He
behoved to come, who should build an house of God, not of wood and stone, but of
men, such as we rejoice He does build. For to this house, that is, to believers,
the apostle saith, ``The temple of God is holy, which temple ye are.''\par
{\smallskip\noindent\small ---St Augustine, City of God, XVII.8\par}
\end{quote}

\noindent St Paul closes his letter to the Romans by praising God for ``the
revelation of the mystery which was kept secret from eternity,'' and in the
Gospel Gabriel tells Mary that the Lord God will give her Son ``the throne of
David his father.'' The promise to David and the consent of Nazareth stand side
by side, and Mary says, ``Behold the handmaid of the Lord: be it done to me
according to thy word.'' The Office of Readings has no second reading of its
own for the Fourth Sunday and takes that of the date; when the Sunday falls on
20 December, it reads St Bernard's homily on Mary's reply
(Section~\ref{sec:last-days}).

In Year A the Gospel turns to Joseph, with Isaiah's sign of the Virgin who shall
conceive. In Year C the prophet Micah names the town:

\begin{quote}
And thou Bethlehem Ephrata, art a little one among the thousands of Juda, out of
thee shall he come forth unto me that is to be the ruler in Israel: and his going
forth is from the beginning, from the days of eternity.\par
{\smallskip\noindent\small ---Micah 5:2 (5:1 in the Lectionary)\par}
\end{quote}

\noindent ``This prophet predicts the very place in which Christ was born,'' wrote
St Augustine; and the Gospel that year is the Visitation. The commentary on the
calendar observed that the readings make this
Sunday, as it were, the Sunday of the Fathers of the Old Testament and of the
Blessed Virgin Mary.

The Collect of the Fourth Sunday is a prayer that many Catholics already know by
heart. Its body is the prayer that closes the Angelus, as Pope St Paul VI
noted (\emph{Marialis cultus}, n.~41 and note). The Missal ends it with the
full conclusion proper to a Collect; the Angelus ends it more briefly:

\medskip
\noindent \latin{Oremus.}\par
\noindent \latin{Gratiam tuam, quaesumus, Domine, mentibus nostris infunde: ut qui, angelo nuntiante, Christi Filii tui Incarnationem cognovimus, per Passionem ejus et Crucem ad resurrectionis gloriam perducamur. Per eundem Christum Dominum nostrum. R. Amen.}\par

\medskip

\noindent Let us pray.\par
\noindent Pour forth, we beseech Thee, O Lord, thy grace into our hearts; that we, to whom the Incarnation of Christ thy Son was made known by the message of an angel, may by his Passion and Cross be brought to the glory of his resurrection. Through the same Christ our Lord. R. Amen.\par
{\smallskip\noindent\small ---The prayer of the Angelus, \emph{Raccolta}\par}
\medskip

\noindent The prayer is old. Its earliest Roman home is the Gregorian
Sacramentary, which gives it at the end of the Mass of the Annunciation, and
before the reform the Missal set it after Communion in the Mass of our Lady in
Advent. It asks in one breath for the whole mystery: the angel's message, the Passion
and the Cross, and the glory of the Resurrection. On the last Sunday of Advent
the Church prays it over the Child not yet born.
